\chapter{Theoretical Analysis and Simulation}
\label{chap:theory-simulation}

This chapter presents the theoretical analysis and numerical simulation of the lateral dynamics of the unicycle robot. The purpose of this analysis is to evaluate the feasibility of the proposed balancing method and controller design before conducting experiments on the physical system.

\section{Modeling of the Lateral Dynamics}

\subsection{Reduced Lateral-Dynamics Model}

The complete unicycle robot contains strongly coupled longitudinal and lateral motions, as illustrated in Figure~\ref{fig:fullnotation}. Developing and analyzing a full model that includes all of these motions is complicated. Therefore, in this chapter, the longitudinal and lateral dynamics are assumed to be approximately decoupled, allowing the lateral motion to be studied independently.

\begin{figure}[htbp]
  \centering
  \includegraphics[width=0.75\textwidth]{02/fullnotation.png}
  \caption{Full mechanical model and notation of the unicycle robot.}
  \label{fig:fullnotation}
\end{figure}

The resulting reduced model of the lateral dynamics is shown in Figure~\ref{fig:partialmodel}. The reduced model is described using the generalized coordinates

\begin{equation}
q_1=\theta,
\qquad
q_2=r-R\theta,
\label{eq:generalized-coordinates}
\end{equation}

where $q_1$ is the lateral lean angle and $q_2$ is the rod position relative to the lateral motion of the wheel center. The corresponding generalized velocities are

\begin{equation}
u_1=\dot{q}_1=\dot{\theta},
\qquad
u_2=\dot{q}_2=\dot{r}-R\dot{\theta}.
\label{eq:generalized-velocities}
\end{equation}

Therefore, the physical rod position and velocity can be recovered from

\begin{equation}
r=q_2+Rq_1,
\qquad
\dot{r}=u_2+Ru_1.
\label{eq:rod-position-recovery}
\end{equation}

To simplify the notation, the constant rotational inertia is defined as

\begin{equation}
J
=
m_wR^2
+
m_b(R+h)^2
+
I_w
+
I_b
+
I_{\mathrm{rod}},
\label{eq:constant-J}
\end{equation}

and the combined gravitational coefficient is defined as

\begin{equation}
G
=
m_wgR
+
m_bg(R+h).
\label{eq:constant-G}
\end{equation}

Here, $m_b$ is the battery mass, which is modeled as the primary pendulum mass of the robot. The parameter $G$ combines the gravitational moments produced by the wheel and the battery assembly.

Using these definitions, the nonlinear lateral dynamics can be written compactly as

\begin{equation}
\begin{bmatrix}
J+m_{\mathrm{rod}}r^2 & 0 \\
0 & m_{\mathrm{rod}}
\end{bmatrix}
\begin{bmatrix}
\dot{u}_1 \\
\dot{u}_2
\end{bmatrix}
=
\begin{bmatrix}
FR
-m_{\mathrm{rod}}gr\cos\theta
+G\sin\theta
-m_{\mathrm{rod}}r
\left(
2u_1u_2+Ru_1^2
\right)
\\[6pt]
F
-m_{\mathrm{rod}}g\sin\theta
+m_{\mathrm{rod}}ru_1^2
\end{bmatrix}.
\label{eq:simplified-lateral-model}
\end{equation}

The four-dimensional state vector is defined as

\begin{equation}
\mathbf{x}
=
\begin{bmatrix}
x_1 \\
x_2 \\
x_3 \\
x_4
\end{bmatrix}
=
\begin{bmatrix}
q_1 \\
q_2 \\
u_1 \\
u_2
\end{bmatrix}
=
\begin{bmatrix}
\theta \\
r-R\theta \\
\dot{\theta} \\
\dot{r}-R\dot{\theta}
\end{bmatrix}.
\label{eq:lateral-state-vector}
\end{equation}

In terms of the state variables, the physical rod position is

\begin{equation}
r=x_2+Rx_1.
\label{eq:rod-position-state}
\end{equation}

The complete four-dimensional model can then be written in the following compact implicit first-order form:

\begin{equation}
\begin{bmatrix}
1 & 0 & 0 & 0 \\
0 & 1 & 0 & 0 \\
0 & 0 & J+m_{\mathrm{rod}}(x_2+Rx_1)^2 & 0 \\
0 & 0 & 0 & m_{\mathrm{rod}}
\end{bmatrix}
\begin{bmatrix}
\dot{x}_1 \\
\dot{x}_2 \\
\dot{x}_3 \\
\dot{x}_4
\end{bmatrix}
=
\begin{bmatrix}
x_3
\\[4pt]
x_4
\\[4pt]
FR
-m_{\mathrm{rod}}g(x_2+Rx_1)\cos x_1
+G\sin x_1 \\
-m_{\mathrm{rod}}(x_2+Rx_1)
\left(
2x_3x_4+Rx_3^2
\right)
\\[6pt]
F
-m_{\mathrm{rod}}g\sin x_1
+m_{\mathrm{rod}}(x_2+Rx_1)x_3^2
\end{bmatrix}.
\label{eq:four-dimensional-lateral-model}
\end{equation}

Equation~\eqref{eq:four-dimensional-lateral-model} is used as the nonlinear plant model for the controller design and numerical simulations presented in the following sections.